\begin{table}[htbp]
\centering
\small
\noindent\textbf{EmoBank 内部認知結合度：問い「他者の感情を評価した変位 $\Delta\text{Reader}$ と、自身の状態として報告した変位 $\Delta\text{Self}$ はモデル内部で連動しているか」}\par\vspace{1ex}
\caption{モデルの他者認識変位と自己報告変位の相関 $\Delta r$（$N=192$ ペア、95\%ブートストラップ信頼区間）。}
\label{tab:behavioral_emobank_coupling}
\begin{tabular}{ll ccc}
\toprule
\textbf{Family} & \textbf{Model Variant} & \textbf{Valence ($\Delta r$ [95\% CI])} & \textbf{Arousal ($\Delta r$ [95\% CI])} & \textbf{Dominance ($\Delta r$ [95\% CI])} \\
\midrule
\multirow{2}{*}{\textbf{Qwen 2.5 1.5B}} & Base       & 0.809 [0.755, 0.854] & 0.866 [0.832, 0.892] & 0.956 [0.944, 0.966] \\
                                 & Instruct   & 0.892 [0.859, 0.918] & 0.917 [0.892, 0.938] & 0.810 [0.748, 0.857] \\
\midrule
\multirow{2}{*}{\textbf{Llama 3.2 1B}} & Base       & 0.662 [0.570, 0.733] & 0.876 [0.839, 0.904] & 0.904 [0.876, 0.926] \\
                                 & Instruct   & 0.756 [0.675, 0.817] & 0.754 [0.697, 0.807] & 0.920 [0.895, 0.941] \\
\midrule
\multirow{2}{*}{\textbf{Gemma 3 1B}} & Base       & 0.813 [0.704, 0.888] & 0.782 [0.685, 0.853] & 0.938 [0.891, 0.964] \\
                                 & Instruct   & 0.621 [0.517, 0.713] & 0.780 [0.660, 0.857] & 0.901 [0.862, 0.931] \\
\midrule
\multirow{2}{*}{\textbf{OLMo 2 1B}} & Base       & 0.559 [0.444, 0.660] & 0.600 [0.469, 0.708] & 0.709 [0.631, 0.767] \\
                                 & Instruct   & 0.851 [0.811, 0.885] & 0.931 [0.908, 0.950] & 0.931 [0.908, 0.949] \\
\bottomrule
\end{tabular}
\vspace{1ex}
\begin{minipage}{\linewidth}
\footnotesize
\textbf{Note:} 全てのファミリーおよびバリアントにおいて、他者認識と自己報告の結合相関 $\Delta r$ は極めて高く（$r > 0.55, p < 10^{-16}$）、モデルが他者の感情を高く評価した刺激に対して自身の自己報告も連動して高く報告する内部的一貫性（認知結合）が確認された。
\end{minipage}
\end{table}